\scanpage{70}
\mspage{76}

\uline{$\mathbb D$-fibrations}

\uline{Prop 1.} Soit $f:X\longrightarrow Y$ dans $\mathrm{Diag}$, Conditions équivalentes :

a) Pour \deleted{\unsure{$y\longrightarrow y'$}} tout $y\in Y$, l'inclusion $X_y\longrightarrow X_{/y}$ est une $\mathbb D$-équivalence, et pour toute flèche $y_1^{\deleted{\unsure{*}}}\longrightarrow y_0$ de $Y$, $X_{/y_1}\longrightarrow X_{/y_0}$ est une $\mathbb D$-équivalence

b) Pour tout \deleted{\unsure{faisceau}} coefficient \deleted{\unsure{loc.}} \uline{constant} $F$ sur $Y$, $f_*f^*F$ est loc. constant\textsuperscript{+}, et ``se calcule fibre par fibre'', l'hom can
\[
f_* f^*(F)_y \unsure{\simeq|\xrightarrow{\sim}} H^{\unsure{\bullet|\cdot|\prime}}_{\mathbb D}(X_y,F\mid X_y).
\]

% slanted
\margin{\textsuperscript{+} Par là on entend que \deleted{\unsure{$\forall$}} $u:y\longrightarrow y'$ dans $Y$, $\unsure{C|G}_{y'}\longrightarrow C_y$ est iso. Sans axiome suppl. sur $\mathbb D$, \deleted{\unsure{$C$}} il n'est pas clair que cela implique que les ``restrictions'' de $C$ aux $Y/y$, ou aux $y\backslash Y$, soient des coefficients constants. Voir cependant l'axiome $\unsure{\mathrm{Der}\,3\,a')|\mathrm{Der}\,3\,a)}$, page \unsure{97|47}, qui assure que c'est pareil.}

Ces conditions sont équivalentes, et \unsure{impliquent|implique} \deleted{que} la validité de b) pour \inserted{tout} $F$ \uline{loc.} constant sur $Y$.

La \inserted{première} condition \inserted{du} a) signifie que \inserted{$\forall y\in Y$, et} $\forall M\in\mathrm{Ob}\,A$, \deleted{\unsure{en tant}}
\[
H^\bullet_{\mathbb D}(X_{/y},M)\longrightarrow H^\bullet_{\mathbb D}(X_y,M)
\]
est un iso. Or cette flèche s'interprète comme
\[
f_*(M_X)_{\unsure{x|y}}\longrightarrow H^\bullet_{\mathbb D}(X_y,M)
\]
est iso, i.e. que $f_*(M_X)$ se calcule fibre par fibre. Comme $M_X\simeq f^*(M_Y)$, c'est la deuxième condition b).

La deuxième condition du a), pour $y_1\longrightarrow y_0$ signifie que pour $M$ dans $A$,
\[
H^\bullet_{\mathbb D}(X_{/y_0},M)\longrightarrow H^\bullet_{\mathbb D}(X_{/y_1},M)
\]
est iso, \deleted{donc que} or cette flèche s'interprète comme
\[
f_*(M_X)_{y_0}\longrightarrow f_*(M_X)_{y_1},
\]
que ce soit \inserted{tjrs} un iso. est la première condition b). Que b) soit valable aussi pour $F$ loc. constant sur $Y$ est immédiat, car c'est une condition de nature locale p.r. à $Y$.
% archive sheet number: 69
